\documentclass[11pt,a4paper]{ctexart}
\usepackage[margin=2.45cm]{geometry}
\usepackage{amsmath,amssymb,amsthm,mathtools,bm,mathrsfs}
\usepackage{xcolor,booktabs,longtable,array,enumitem,hyperref}
\hypersetup{colorlinks=true,linkcolor=blue!55!black,urlcolor=blue!55!black,citecolor=blue!55!black}
\setlist{nosep}

\newtheorem{theorem}{定理}[section]
\newtheorem{lemma}[theorem]{引理}
\newtheorem{proposition}[theorem]{命题}
\theoremstyle{definition}
\newtheorem{definition}[theorem]{定义}
\newtheorem{remark}[theorem]{说明}

\newcommand{\F}{\mathbb F}
\newcommand{\G}{\F_p^\times}
\newcommand{\what}{\widehat}
\newcommand{\OPEN}{\textcolor{red!70!black}{\textsf{OPEN}}}
\newcommand{\PASS}{\textcolor{green!45!black}{\textsf{PASS}}}
\newcommand{\FAIL}{\textcolor{red!70!black}{\textsf{FAIL}}}
\newcommand{\REDUCED}{\textcolor{orange!80!black}{\textsf{REDUCED}}}

\title{G2 最终封口为何失败，以及应怎样修复\\
\large 从 WJCTLS replay 到 DCCS--KPA8 的逐层审计}
\author{依据最终对话分支重建的技术报告}
\date{2026年8月28日}

\begin{document}
\maketitle

\begin{abstract}
本文说明 balanced G2 为何在多次“即将关闭”之后仍不能通过严格审计。
最早的最终 replay 中，标量乘法 $Q^4/p^5\asymp p^{-1}$ 本身正确，
但 $D,D'$ Mellin inversion 的共轭方向被混用；正确支撑是 reciprocal/projective
fibre，而非 affine fibre，因此后续 additive Moyal/SWAFM 与 WJCTLS 逻辑断开。
后续工作把正确 residual 依次压到 CRAF、centered KPA8 与 diagonal
sum--product charge，但最新链仍缺 DCCS 的 exact operator factorization，
以及同一 normalization 下把实际 $m$-$L^8$ 资源送入 Schatten-$8$ target 的证明。
本文给出可检验的修复顺序、禁止捷径及 Lean proof obligations。
\end{abstract}

\section{最终判决}

截至所给日志末端，最稳妥的状态是
\[
\boxed{\mathrm{G2}\quad\textbf{尚未关闭}.}
\]
更准确地，存在两次不同层级的审计：
\begin{enumerate}
\item 旧 WJCTLS $\to$ affine diagonalization $\to$ SWAFM 链在 Mellin
共轭方向处 \FAIL；
\item 修正后的 reciprocal/projective 链已到
\[
\mathrm{CRAF}\to\mathrm{KPA8}\to\mathrm{DCCS},
\]
但 DCCS 仍只有候选 factorization，没有逐指标证明；KPA8 的数值界也不能由
NOFA 的 $L^2$ Bessel 性质自动推出。
\end{enumerate}
因此“只差一个 $p$ 的 normalization”已经不是当前正确描述；当前缺口是
\emph{operator type splice + normalized high-moment insertion}。

\section{旧 replay：算术正确，逻辑拼接错误}

固定奇素数 $p$，令 $Q=p-1$。日志中的 fully-cross WJCTLS quadratic form为
\[
\mathcal Q
=Q^{-3}\sum_{\gamma,z,z',D,D'}
\widetilde W_\gamma(z,D)
\overline{\widetilde W_\gamma(z',D')}
\mathfrak K(z/z';D,D').
\tag{2.1}
\]
generic Jacobi kernel写成
\[
\mathfrak K(r;D,D')=Q\sqrt p\,\Theta_r(D,D'),
\]
故首个 exact scalar 为
\[
Q^{-3}Q\sqrt p=\frac{\sqrt p}{Q^2}.
\tag{2.2}
\]

后续日志曾把四段 prefactor 乘为
\[
\frac{\sqrt p}{Q^2}\cdot\frac Q{\sqrt p}\cdot pQ
\cdot\left(\frac Q{p^{3/2}}\right)^4
=\boxed{\frac{Q^4}{p^5}}
=p^{-1}\bigl(1+O(p^{-1})\bigr).
\tag{2.3}
\]
这个算术等式是正确的。错误是第三、第四个因子来自一条并不适用于正确 fibre
的 affine-Moyal 链；所以 (2.3) 不是 WJCTLS 的证明证书。

\section{首个致命点：Mellin inverse 的方向}

取单位化 Mellin 约定
\[
\mathcal W_\gamma(z,y)
=Q^{-1/2}\sum_{D\in\widehat\G}\overline{D(y)}
\widetilde W_\gamma(z,D).
\tag{3.1}
\]
由字符正交，其逆变换必须是
\[
\widetilde W_\gamma(z,D)
=Q^{-1/2}\sum_{y\in\G}D(y)\mathcal W_\gamma(z,y).
\tag{3.2}
\]
不能在 (3.1) 使用 $\overline{D(y)}$，却在下游按另一个 convention
使用 inverse。

令
\[
J(D,\overline{D'})=\sum_{x\in\F_p}D(x)\overline{D'(1-x)},
\quad
a=\frac{r-1}{r},\quad b=1-r.
\]
把 (3.2) 代入含 $D(a)D(x)$ 与 $\overline{D'(b(1-x))}$ 的核。
字符正交给出的支撑为
\[
\boxed{yax=1},\qquad \boxed{y'b(1-x)=1},
\tag{3.3}
\]
而不是 $y=ax$、$y'=b(1-x)$。

\begin{proposition}[正确的 reciprocal fibre]
设 $z,z'\in\G$ 且 $z\ne z'$，$r=z/z'$。在 (3.3) 上有
\[
\boxed{z(y^{-1}-1)=z'(y'^{-1}-1).}
\tag{3.4}
\]
\end{proposition}

\begin{proof}
由 $a=(z-z')/z$、$b=(z'-z)/z'$，(3.3) 给
\[
x=\frac{z}{y(z-z')},\qquad
1-x=-\frac{z'}{y'(z-z')}.
\]
相加并使用 $x+(1-x)=1$，得到
$z/y-z'/y'=z-z'$，即 (3.4)。
\end{proof}

旧链使用的是
\(z(y-1)=z'(y'-1)\)。这两个式子不能与同一个 coefficient formula
同时成立；把 $y$ 改名成 $y^{-1}$ 会同步改变 short coefficient 中
$u(yz)$ 的方向，因而不是无害重标号。

\section{为什么 affine SWAFM 不能再接回去}

令
\[
q=z(y^{-1}-1),\qquad w=yz.
\]
则
\[
q=z\left(\frac zw-1\right),\qquad
\boxed{w=\frac{z^2}{z+q}},
\tag{4.1}
\]
而不是 $w=z+q$。故真实 residual 是 projective correlation：
\[
S_\gamma^{\rm rec}(q)
=\sum_\alpha m(\alpha)\overline{m(\alpha\gamma)}
\sum_{z\ne0,-q}\overline{U_\alpha(z)}
U_{\alpha\gamma}\!\left(\frac{z^2}{z+q}\right).
\tag{4.2}
\]
所需目标为
\[
\boxed{
\sum_{\gamma,q}|S_\gamma^{\rm rec}(q)|^2
\ll p^{5/2+o(1)}.}
\tag{4.3}
\]
旧 affine SWAFM 即便作为独立命题成立，也不估计 (4.2)。

\section{正确路线已推进到哪里}

日志随后没有再次强行使用错误 fibre，而是做了如下合法的结构压缩：
\[
\text{reciprocal WSRE}
\Longrightarrow \text{two-chart projective/M\"obius frame}
\Longrightarrow \mathrm{CRAF}
\Longrightarrow \mathrm{KPA8}.
\]
其中最稳定的代数结论是：对 $r\in\F_p\setminus\{0,1\}$，令
\[
x=r^{-1},\qquad z=(1-r)^{-1}.
\]
则
\[
R:=x+z=\frac1{r(1-r)}=xz=:P.
\tag{5.1}
\]
因此 projective Kloosterman curve是 sum--product joint-charge space中的
对角 $R=P$。

定义 exact sum--product charge transform
\[
(\mathcal AF)(R,P)
=\sum_{\substack{x,z\in\G\\x+z=R,\ xz=P}}F(x,z),
\quad R\in\F_p,\ P\in\G,
\tag{5.2}
\]
以及它在 $P$ 坐标上的 unitary character-coordinate 版本
\[
(\mathcal JF)(R,\eta)
=Q^{-1/2}\sum_{\substack{x,z\in\G\\x+z=R}}
F(x,z)\eta(xz).
\tag{5.3}
\]
两者只相差 $P$ 坐标上的 unitary Mellin 变换。直接用字符正交可证
\[
\boxed{\mathcal A^*\mathcal A
=\mathcal J^*\mathcal J=I+\mathsf S-\mathsf D},
\qquad \|\mathcal A\|_{2\to2}=\|\mathcal J\|_{2\to2}\le\sqrt2,
\tag{5.4}
\]
其中 $\mathsf S$ 交换 $(x,z)$，$\mathsf D$ 投影到 $x=z$。
限制 $R=P$ 是 $\mathcal A$ 的 exact-charge target space上的正交投影，范数为 $1$。
在 $\mathcal J$ 的 $(R,\eta)$ 坐标中不能直接比较 $R$ 与角色 $\eta$；必须先
作 inverse Mellin。这一层类型转换是 DCCS 必须显式包含的。

\paragraph{但这仍不等于 KPA8。}
KPA8 的 centered coefficient为
\[
B^\circ(A,r)
=\sum_w y(rw+Ar(1-r))\overline{y(w)}-\mu_y.
\tag{5.5}
\]
用 (5.1) 后，平移量含有
\[
Ar(1-r)=AP^{-1}.
\tag{5.6}
\]
所以 coefficient 不是先固定一个 $F(x,z)$ 再做 $\mathcal J$；它是
依赖 joint product charge $P$ 的 family。正是这一点阻止了
“NOFA $\Rightarrow$ KPA8”的一句话证明。

\section{最新最小接口：DCCS}

令
\[
\mathscr H=\bigoplus_{R\in\G}\ell^2(\F_p),
\]
并在每个 $R$-fibre 定义 additive dilation
\[
(\mathcal U_{\rm fib}f)(R,A)=f(R,A/R).
\tag{6.1}
\]
因为 $A\mapsto A/R$ 是 $\F_p$ 的置换，$\mathcal U_{\rm fib}$ 是酉算子。
令 $\Pi_\Delta$ 为 $R=P$ 的正交投影，$\Pi_0^\perp$ 为删除 additive
zero mode 的正交投影。候选 DCCS identity 是
\[
\boxed{
\mathcal S_{\rm KPA8}
=\Pi_0^\perp\mathcal U_{\rm fib}\Pi_\Delta
\mathcal A\mathcal C_y,}
\tag{6.2}
\]
允许一个显式的 swap symmetrizer、已声明的标量常数，以及当输入最初写在
$(R,\eta)$ 坐标时所需的 inverse Mellin。

\paragraph{当前状态。}
(6.1) 的酉性已证；$\Pi_\Delta$、$\Pi_0^\perp$ 的 contraction 性已证；
但 (6.2) 尚未从 (5.5) 对每个 basis vector逐项推出，故
\[
\boxed{\mathrm{DCCS}\;=\;\OPEN.}
\]

\section{为什么 DCCS 成立后仍需一次高矩回代}

KPA8 的数值目标是
\[
\boxed{\|B^\circ\|_{S_8}^8\ll p^{6+o(1)}.}
\tag{7.1}
\]
若 (6.2) 成立，ideal property 只给“assembly 不再产生额外 rank loss”：
\[
\|\mathcal S_{\rm KPA8}\|_{S_8}
\le \sqrt2\,\|\mathcal C_y\|_{S_8}.
\tag{7.2}
\]
它不会凭空证明右侧达到 (7.1) 的尺度。

日志声称已有 normalized multiplier $m$ 的资源
\[
\sum_\chi|m(\chi)|^4\ll p^{o(1)},
\qquad
\sum_\chi|m(\chi)|^8\ll p^{-1/4+o(1)}.
\tag{7.3}
\]
在最终证明中必须同时给出：$m$ 的精确定义、principal character 的处理、
support、Fourier normalization，以及 (7.3) 如何在 (6.2) 的同一 Hilbert
measure 中成为 (7.1)。否则“已有高矩”仍只是接口名。

\section{另一个不可忽略的外部类型问题：Kable 的 measure}

Kable 的定理确实把 Legendre 与 Soto--Andrade sums 解释为两个 Steinberg
表示的 parabolic Clebsch--Gordan coefficients，并证明完备正交性；但原定理的
Hilbert space 是 $\ell^2(\F_q,m)$，其中 $\pm1$ 的质量为 $q+1$，其他点质量为 $1$。
所以正式拼接必须证明下列三者之一：
\begin{enumerate}
\item 我们的 operator 本来就在同一个 weighted measure 上；
\item 端点权通过显式 diagonal conjugation 被运输；
\item $\pm1$ 分支被精确拆为有限秩 carrier，并在目标尺度内单独估计。
\end{enumerate}
仅写“CG basis 是 unitary”而不说明 measure，不能通过同行或 Lean 审计。

\section{推荐的封口顺序}

\begin{enumerate}[label=\textbf{Step \arabic*.}]
\item \textbf{冻结 convention。} 全文只用一个 additive Fourier、一个
multiplicative Mellin 约定；把每个 inverse 写在定义旁。
\item \textbf{定义带 measure 的有限 Hilbert spaces。} 对每个 index 指明
$\F_p$ 或 $\F_p^\times$、counting 或 probability measure，以及矩阵的 row/column 类型。
\item \textbf{先证 DCCS kernel identity。} 对输入 basis
$\delta_{x,z,A}$ 逐项计算 (6.2) 两边；不用 $\asymp$、不用范数、也不用外部估计。
\item \textbf{处理 $R=0$、$r=0,1$、$x=z$、$\pm1$。} 每个 exceptional
branch须有显式投影或有限秩项，不能藏入 $p^{o(1)}$。
\item \textbf{证明 Schatten transport。} 从 exact identity 用
$\|AXB\|_{S_8}\le\|A\|_{op}\|X\|_{S_8}\|B\|_{op}$，记录唯一的 $\sqrt2$。
\item \textbf{把 (7.3) 送入 actual short tensor。} 这一步必须是等式或一次
明确 Hölder，不得把 $L^2$ Bessel 当作 $S_8$ saving。
\item \textbf{做一次从 KPA8 到 G2 的全标量 replay。} 顺序固定为
\[
\mathrm{DCCS}\Rightarrow\mathrm{KPA8}\Rightarrow\mathrm{CRAF}
\Rightarrow\mathrm{WSRE}\Rightarrow\mathrm{WJCTLS}\Rightarrow\mathrm{G2},
\]
每个箭头给 exact scalar 与 domain bijection。
\end{enumerate}

\paragraph{如果 Step 3 失败。}
不要再引入新命名。把两边相减，得到显式 commutator/residual
\[
\mathcal R_{\rm DCCS}
=\mathcal S_{\rm KPA8}
-\Pi_0^\perp\mathcal U_{\rm fib}\Pi_\Delta\mathcal A\mathcal C_y,
\]
然后只估计这个 residual。这样能判断问题究竟是一个有限秩 endpoint、
一个 swap multiplicity，还是 $P$-dependent dilation 与 joint transform 不交换。

\section{禁止再次使用的捷径}

\begin{longtable}{p{0.39\textwidth}p{0.50\textwidth}}
\toprule
禁止陈述 & 原因\\
\midrule
\endfirsthead
\toprule
禁止陈述 & 原因\\
\midrule
\endhead
$Q^4/p^5\sim p^{-1}$ 因而 WJCTLS 已证
& 标量正确，但中间 affine-Moyal 箭头不适用于 reciprocal fibre。\\
$y\mapsto y^{-1}$ 是无害重命名
& 它同时改变 coefficient $u(yz)$ 的方向；混用两种 convention 是 type error。\\
entrywise bijection/shear 保持 Schatten 范数
& 只有左右酉乘与合适的 unitary conjugation 保持 singular values；一般 entry
重排不保持。\\
NOFA $L^2$ Bessel 自动给 KPA8 $S_8$
& Bessel 只控制 assembly multiplicity，不创造高矩 saving。\\
multiplicity-one 自动使 coefficient operator diagonal
& 还需证明实际 operator 与 irreducible direct sum 的 intertwining identity。\\
Kable 正交基在 counting measure 下无条件酉
& 原文有特殊端点权；measure 必须运输。\\
数值小素数测试证明渐近定理
& 数值测试只能发现错误，不能提供一致常数与所有 $p$ 的证明。\\
\bottomrule
\end{longtable}

\section{Lean theorem ledger}

\begin{center}
\begin{tabular}{p{0.43\textwidth}p{0.18\textwidth}p{0.26\textwidth}}
\toprule
obligation & 状态 & 形式化备注\\
\midrule
Mellin inversion (3.1)--(3.2) & \PASS & 有限字符正交\\
reciprocal fibre (3.4) & \PASS & 有限域代数，分母 guard\\
projective map (4.1) & \PASS & 双射及 $z+q\ne0$\\
joint Gram (5.4) & \PASS & 有限和，Vi\`ete fibres\\
$R=P$ projection & \PASS & 正交投影\\
fibre dilation (6.1) & \PASS & fibrewise permutation\\
DCCS exact factorization (6.2) & \OPEN & 必须逐 basis 证明\\
Kable measure splice & \OPEN & weighted/unweighted 类型\\
$m$-$L^8$ 到 actual tensor & \OPEN & exact normalization\\
KPA8 numerical target & \OPEN & 依赖前两项\\
WJCTLS/G2 final replay & \OPEN & 只能在 KPA8 后执行\\
\bottomrule
\end{tabular}
\end{center}

\section{最简研究建议}

下一刀应只做 DCCS：把 (5.5) 在 $(x,z,A)$ basis 上展开，与 (6.2) 两边逐项比较。
这一步是有限维线性代数，不需要新的 Burgess、Weil、Katz 或 non-abelian
amplification。只有 exact identity 通过后，才值得消费 (7.3) 并跑最后 exponent ledger。
若 DCCS 通过，G2 很可能确实只剩一次机械 normalization；若不通过，
残差本身将成为唯一且诚实的新目标。

\begin{thebibliography}{9}
\bibitem{Kable2002}
A.~C. Kable,
\emph{Legendre sums, Soto--Andrade sums and Kloosterman sums},
Pacific J. Math. 206 (2002), 139--157.
\url{https://msp.org/pjm/2002/206-1/pjm-v206-n1-p09-s.pdf}

\bibitem{Pascadi2026}
A. Pascadi,
\emph{Non-abelian amplification and bilinear forms with Kloosterman sums},
arXiv:2511.08445v2 (2026), especially Proposition 5.1.
\url{https://arxiv.org/abs/2511.08445}

\bibitem{BP2026}
V. Blomer and A. Pascadi,
\emph{Bilinear forms with Kloosterman sums via quadratic characters},
arXiv:2607.24311v1 (2026).
\url{https://arxiv.org/abs/2607.24311}
\end{thebibliography}

\end{document}
